\documentclass[UTF8,12pt,a4paper,fontset=none]{ctexart}
\usepackage[a4paper,left=1.82cm,right=1.82cm,top=1.72cm,bottom=1.82cm,headheight=15pt]{geometry}
\usepackage{fontspec,xeCJK}
\setmainfont{Liberation Serif}
\setsansfont{DejaVu Sans}
\setmonofont{DejaVu Sans Mono}
\setCJKmainfont[BoldFont=Noto Serif CJK SC Bold]{Noto Serif CJK SC}
\setCJKsansfont[BoldFont=Noto Sans CJK SC Bold]{Noto Sans CJK SC}
\setCJKmonofont{Noto Sans Mono CJK SC}
\usepackage{amsmath,amssymb,mathtools,bm}
\usepackage{booktabs,longtable,tabularx,array,multirow}
\usepackage{enumitem,xcolor}
\usepackage[most]{tcolorbox}
\usepackage{fancyhdr,titlesec,hyperref,cleveref,lastpage}
\usepackage{tikz,float,caption,listings}
\usetikzlibrary{arrows.meta,positioning,calc,fit,backgrounds,shapes.geometric}
\definecolor{mainblue}{RGB}{39,82,138}
\definecolor{deepblue}{RGB}{25,55,95}
\definecolor{softblue}{RGB}{235,243,252}
\definecolor{softgreen}{RGB}{238,248,241}
\definecolor{softorange}{RGB}{255,246,232}
\definecolor{softgray}{RGB}{245,246,248}
\definecolor{warnred}{RGB}{160,48,48}
\hypersetup{unicode=true,colorlinks=true,linkcolor=deepblue,urlcolor=mainblue,citecolor=deepblue}
\setlength{\parindent}{2em}
\setlength{\parskip}{0.36em}
\setlength{\tabcolsep}{5pt}
\renewcommand{\arraystretch}{1.2}
\allowdisplaybreaks[4]
\emergencystretch=3em
\sloppy
\pagestyle{fancy}\fancyhf{}\fancyfoot[C]{\small 第 \thepage\ 页，共 \pageref{LastPage} 页}\renewcommand{\headrulewidth}{0.4pt}
\titleformat{\section}{\Large\bfseries\color{deepblue}}{\thesection}{0.65em}{}
\titleformat{\subsection}{\large\bfseries\color{mainblue}}{\thesubsection}{0.55em}{}
\titleformat{\subsubsection}{\normalsize\bfseries}{\thesubsubsection}{0.5em}{}
\numberwithin{equation}{section}
\newcommand{\R}{\mathbb{R}}
\newcommand{\E}{\mathbb{E}}
\newcommand{\Prob}{\mathbb{P}}
\newcommand{\Loss}{\mathcal{L}}
\newcommand{\norm}[1]{\left\lVert #1\right\rVert}
\newcommand{\ip}[2]{\left\langle #1,#2\right\rangle}
\newcommand{\tr}{\operatorname{Tr}}
\newcommand{\diag}{\operatorname{diag}}
\newcommand{\rank}{\operatorname{rank}}
\newcommand{\softmax}{\operatorname{softmax}}
\newcommand{\argmin}{\operatorname*{arg\,min}}
\newcommand{\argmax}{\operatorname*{arg\,max}}
\newcommand{\sg}{\operatorname{sg}}
\newcommand{\KL}{D_{\mathrm{KL}}}
\newcommand{\dd}{\mathrm{d}}
\newtcolorbox{keybox}[1][]{enhanced,breakable,colback=softblue,colframe=mainblue,boxrule=0.8pt,arc=1.5mm,left=2mm,right=2mm,top=1.2mm,bottom=1.2mm,title={#1},fonttitle=\bfseries}
\newtcolorbox{examplebox}[1][]{enhanced,breakable,colback=softgreen,colframe=green!45!black,boxrule=0.7pt,arc=1.5mm,left=2mm,right=2mm,top=1.2mm,bottom=1.2mm,title={#1},fonttitle=\bfseries}
\newtcolorbox{notebox}[1][]{enhanced,breakable,colback=softorange,colframe=orange!70!black,boxrule=0.7pt,arc=1.5mm,left=2mm,right=2mm,top=1.2mm,bottom=1.2mm,title={#1},fonttitle=\bfseries}
\newtcolorbox{criticalbox}[1][]{enhanced,breakable,colback=red!3,colframe=warnred,boxrule=0.8pt,arc=1.5mm,left=2mm,right=2mm,top=1.2mm,bottom=1.2mm,title={#1},fonttitle=\bfseries}
\lstset{basicstyle=\ttfamily\small,breaklines=true,frame=single,backgroundcolor=\color{softgray},numbers=left,numberstyle=\tiny,showstringspaces=false,columns=fullflexible}
\hypersetup{pdftitle={39 Continuous Thought Machines 数学过程详解},pdfauthor={OpenAI}}
\fancyhead[L]{\small 39 Continuous Thought Machines}
\fancyhead[R]{\small 数学过程详解}
\setlength{\abovedisplayskip}{0.82em}
\setlength{\belowdisplayskip}{0.82em}
\setlength{\abovedisplayshortskip}{0.45em}
\setlength{\belowdisplayshortskip}{0.45em}
\begin{document}
\begin{center}
{\LARGE\bfseries 39\_Continuous Thought Machines 数学过程详解}\par
\vspace{0.35em}
{\large 神经元历史模型、同步矩阵与确定度自适应计算}
\end{center}
\vspace{0.45em}
\begin{keybox}[本文只处理真正需要推导的部分]
本文逐式推导内部 tick、私人 NLM、同步 Gram 矩阵、指数衰减递推、pair subsampling、同步注意力、熵确定度、双 tick 损失和停机期望。全文按“公式从哪里来--每个符号是什么--中间步骤为何成立--维度和复杂度如何核对--论文结论依赖哪些假设”的顺序展开；不复述实验表格，也不使用与论文无关的通用背景凑页数。
\end{keybox}
\tableofcontents
\newpage

\section{CTM 的内部时间轴}
静态输入也在内部 tick $t=1,\ldots,T$ 上迭代。神经元状态
\begin{equation}
 z_t\in\R^D,
\end{equation}
与数据注意力输出 $o_t$ 经过共享突触网络
\begin{equation}
 \boxed{a_t=f_{\theta_{\rm syn}}([z_t;o_t])\in\R^D}.
\end{equation}
最近 $M$ 个 pre-activation 组成
\begin{equation}
 A_t=[a_{t-M+1},\ldots,a_t]\in\R^{D\times M}.
\end{equation}

\section{每个神经元的私人时间模型}
第 $d$ 个神经元取历史行向量 $A_t^{(d)}\in\R^M$，用独立参数
\begin{equation}
 z_{t+1,d}=g_{\theta_d}(A_t^{(d)}).
\end{equation}
若一层 MLP 宽度 $h$：
\begin{equation}
 g_{\theta_d}(u)=w_{2,d}^\top\sigma(W_{1,d}u+b_{1,d})+b_{2,d}.
\end{equation}
参数量约
\begin{equation}
 P_{\rm NLM}=D(hM+h+h+1)=O(DhM),
\end{equation}
区别于标准 RNN 的共享时间更新矩阵：CTM 的每个神经元对自己的历史有不同滤波器。

\section{同步矩阵作为 latent representation}
收集后激活历史
\begin{equation}
 Z_t=[z_1,z_2,\ldots,z_t]\in\R^{D\times t}.
\end{equation}
同步矩阵
\begin{equation}
 \boxed{S_t=Z_tZ_t^\top\in\R^{D\times D}},
\end{equation}
元素
\begin{equation}
 S_{t,ij}=\sum_{\tau=1}^{t}z_{\tau,i}z_{\tau,j}.
\end{equation}
它是未中心化时间相关矩阵。对角线是神经元历史能量，非对角线是共同激活；比单帧 $z_t$ 包含更长时间结构。

\section{归一化指数衰减同步}
每对神经元有 $r_{ij}\ge0$，权重
\begin{equation}
 w_\tau=e^{-r_{ij}(t-\tau)}.
\end{equation}
论文同步量
\begin{equation}
 \boxed{
 S_{t,ij}^{(r)}=
 \frac{\sum_{\tau=1}^{t}w_\tau z_{\tau,i}z_{\tau,j}}
 {\sqrt{\sum_{\tau=1}^{t}w_\tau}}}.
\end{equation}
$r=0$ 时平等使用全历史；$r$ 大时只看最近 tick。有效时间常数约
\begin{equation}
 \tau_{\rm eff}\approx1/r
\end{equation}
（离散小 $r$ 情形）。

\section{同步的递推计算}
定义分子和归一化量
\begin{equation}
 P_t=\sum_{\tau=1}^{t}e^{-r(t-\tau)}z_{\tau,i}z_{\tau,j},
 \quad Q_t=\sum_{\tau=1}^{t}e^{-r(t-\tau)}.
\end{equation}
则
\begin{align}
 P_{t+1}&=e^{-r}P_t+z_{t+1,i}z_{t+1,j},\\
 Q_{t+1}&=e^{-r}Q_t+1.
\end{align}
同步量 $S_{t,ij}=P_t/\sqrt{Q_t}$。每个选中 neuron pair 每 tick 只需 $O(1)$ 更新，无需保存完整 $Z_t$。

\section{Pair Subsampling 的维度}
完整 $S_t$ 有 $D^2$ 元素。随机选 $D_{\rm action}$ 对形成
\begin{equation}
 s_t^{\rm action}\in\R^{D_{\rm action}},
\end{equation}
再投影为注意力 query
\begin{equation}
 q_t=W_{\rm in}s_t^{\rm action}.
\end{equation}
另选 $D_{\rm out}$ 对
\begin{equation}
 y_t=W_{\rm out}s_t^{\rm out}.
\end{equation}
计算从 $O(D^2t)$ 降为 $O((D_{\rm action}+D_{\rm out})t)$，但随机子采样只保留部分二阶关系。

\section{同步如何驱动外部注意力}
数据特征 $K,V=F(x)$，query 来自内部同步：
\begin{equation}
 o_t=\operatorname{Attention}(q_t,K,V)
 =\softmax\left(\frac{q_tK^\top}{\sqrt d}\right)V.
\end{equation}
因此下一个内部状态依赖
\begin{equation}
 z_{t+1}=G(z_{\le t},\operatorname{Attention}(S_t,F(x))).
\end{equation}
注意位置不是固定扫描，而是由过去神经动态的二阶统计自适应决定。

\section{确定度的熵定义}
类别分布 $p_t=\softmax(y_t)$，熵
\begin{equation}
 H_t=-\sum_{c=1}^{C}p_{t,c}\log p_{t,c}.
\end{equation}
最大熵为 $\log C$，归一化确定度
\begin{equation}
 \boxed{C_t=1-\frac{H_t}{\log C}}.
\end{equation}
均匀分布时 $C_t=0$，one-hot 极限时 $C_t=1$。这衡量分布尖锐度，不保证预测正确。

\section{双 tick 损失}
每 tick 交叉熵
\begin{equation}
 L_t=-\log p_t(y).
\end{equation}
选择
\begin{equation}
 t_1=\argmin_tL_t,
 \qquad t_2=\argmax_tC_t,
\end{equation}
总损失
\begin{equation}
 \boxed{\Loss=\frac{L_{t_1}+L_{t_2}}{2}}.
\end{equation}
第一项奖励任意时刻出现正确解；第二项要求“最自信的时刻也要正确”，抑制错误高置信度。

\section{argmin/argmax 选择的梯度}
选择索引视为离散常数，反向只通过两个 tick：
\begin{equation}
 \nabla_\theta\Loss=\frac12
 (\nabla_\theta L_{t_1}+\nabla_\theta L_{t_2}).
\end{equation}
索引切换边界处损失不可微，但几乎处处可用子梯度。缺点是其它 tick 无直接监督；共享递归参数仍会通过 BPTT 间接收到梯度。

\section{自适应停机}
给定阈值 $\tau$，停机时间
\begin{equation}
 T_\tau=\min\{t:C_t\ge\tau\}\wedge T.
\end{equation}
期望计算
\begin{equation}
 \E[T_\tau]=\sum_{t=0}^{T-1}\Pr(T_\tau>t).
\end{equation}
提高 $\tau$ 通常增加计算和准确率；若确定度未校准，早停可能把自信错误提前输出。

\section{与普通 RNN 的状态维度差异}
普通 RNN 用 $z_t\in\R^D$ 作为全部状态。CTM 用同步向量，潜在二阶维度可达
\begin{equation}
 \frac{D(D+1)}2.
\end{equation}
即使只采样 $P$ 对，仍等价于从高维二阶特征映射
\begin{equation}
 \varphi(Z_t)=\{\ip{Z_{t,i}}{Z_{t,j}}\}_{(i,j)}
\end{equation}
中选坐标。这类似二次核特征，能线性读出单个神经元状态无法表示的相位/同步关系。

\section{同步尺度不变性与局限}
若所有历史状态乘 $a$，
\begin{equation}
 S_t\mapsto a^2S_t.
\end{equation}
因此同步不是幅值不变。可选的相关系数归一化
\begin{equation}
 \rho_{ij}=\frac{S_{ij}}{\sqrt{S_{ii}S_{jj}}}
\end{equation}
会去除尺度，但论文保留能量信息。其代价是激活幅值可能主导同步，需依赖归一化和训练稳定性。
% AUGMENTED_DETAILED_MATH_2

\section{私人神经元时间模型的卷积形式}
CTM 每个神经元 $i$ 不只看当前 pre-activation，而是看长度 $M$ 的历史
\begin{equation}
 a_{t-M+1:t,i}.
\end{equation}
线性私人模型可写成一维因果卷积
\begin{equation}
 z_{t,i}=\sum_{m=0}^{M-1}w_{i,m}a_{t-m,i}+b_i.
\end{equation}
不同神经元拥有不同核 $w_i$，因此同一输入脉冲可产生不同延迟、振荡或积分响应。若再接非线性 $\phi$，
\begin{equation}
 z_{t,i}=\phi(w_i^\top a_{t-M+1:t,i}+b_i).
\end{equation}
这相当于每个神经元有私人 FIR 动力系统，而普通 MLP 神经元只对应 $M=1$。

\section{时间核的频率响应}
线性核的离散 Fourier 响应
\begin{equation}
 H_i(e^{i\omega})=\sum_{m=0}^{M-1}w_{i,m}e^{-i\omega m}.
\end{equation}
不同 $w_i$ 可形成低通、高通或带通神经元。同步表示随后计算不同频率响应之间的相关性，因此 CTM 的时间处理不只是“多迭代几次”，而是在神经元层面学习不同时间滤波器。

\section{同步矩阵是 Gram 矩阵}
收集每个神经元历史向量
\begin{equation}
 Z_i=[z_{1,i},\ldots,z_{t,i}]^\top.
\end{equation}
同步
\begin{equation}
 S_{ij}=Z_i^\top Z_j.
\end{equation}
矩阵形式
\begin{equation}
 S=Z^\top Z.
\end{equation}
因此
\begin{equation}
 S\succeq0,
 \qquad \rank(S)\le t.
\end{equation}
对任意 $c$，
\begin{equation}
 c^\top Sc=\norm{Zc}^2\ge0.
\end{equation}
同步表示本质是神经元时间轨迹的二阶 Gram 特征。

\section{指数衰减同步的归一化}
权重
\begin{equation}
 w_\tau=e^{-r(t-\tau)}.
\end{equation}
加权内积
\begin{equation}
 P_{ij}=\sum_\tau w_\tau z_{\tau,i}z_{\tau,j}.
\end{equation}
论文除以 $\sqrt{\sum w_\tau}$ 而不是 $\sum w_\tau$：
\begin{equation}
 S_{ij}=P_{ij}/\sqrt{Q},
 \qquad Q=\sum w_\tau.
\end{equation}
若 $z_i,z_j$ 是独立零均值、方差固定随机序列，则 $P_{ij}$ 标准差约与 $\sqrt{\sum w_\tau^2}$ 成正比；平方根归一化使统计尺度随有效窗口增长更慢，接近保持随机相关量的方差，而非计算普通加权平均。

\section{递推与时间复杂度}
令衰减因子 $\lambda=e^{-r}$，
\begin{equation}
 P_t=\sum_{\tau=1}^{t}\lambda^{t-\tau}z_{\tau,i}z_{\tau,j}.
\end{equation}
则
\begin{equation}
 P_{t+1}=\lambda P_t+z_{t+1,i}z_{t+1,j},
\end{equation}
\begin{equation}
 Q_{t+1}=\lambda Q_t+1.
\end{equation}
若只采样 $P$ 对神经元，每 tick 成本 $O(P)$、状态 $O(P)$；若完整计算所有对，则成本和存储 $O(D^2)$。

\section{随机 pair subsampling 的估计性质}
完整同步特征有 $D(D+1)/2$ 个。若均匀随机采样集合 $\mathcal P$，大小 $P$，对某线性读出
\begin{equation}
 y=\sum_{i\le j}a_{ij}S_{ij}
\end{equation}
可构造无偏估计
\begin{equation}
 \widehat y=\frac{D(D+1)/2}{P}
 \sum_{(i,j)\in\mathcal P}a_{ij}S_{ij}.
\end{equation}
方差随 $1/P$ 下降，但论文实际把采样坐标输入可训练投影，并不显式做该重要性缩放；它学习的是一个随机二阶子空间，而非完整同步读出的无偏 Monte Carlo 近似。

\section{同步注意力的维度链}
选中同步向量
\begin{equation}
 s_t\in\R^P,
\end{equation}
线性投影
\begin{equation}
 q_t=W_qs_t,
 \qquad W_q\in\R^{d_k\times P}.
\end{equation}
外部特征 $K\in\R^{N\times d_k}$、$V\in\R^{N\times d_v}$，注意力
\begin{equation}
 o_t=\softmax(Kq_t/\sqrt{d_k})^\top V.
\end{equation}
因此外部读取位置由内部二阶历史 $s_t$ 决定，梯度链为
\begin{equation}
 \frac{\partial\Loss}{\partial s_t}
 =W_q^\top
 \frac{\partial q_t}{\partial s_t}^\top
 \frac{\partial o_t}{\partial q_t}^\top
 \frac{\partial\Loss}{\partial o_t}.
\end{equation}
同步不是仅用于最终分类，而会反向塑造每个 tick 的观察策略。

\section{熵确定度的边界与梯度}
归一化确定度
\begin{equation}
 C(p)=1+\frac1{\log K}\sum_{k=1}^{K}p_k\log p_k.
\end{equation}
由熵界 $0\le H(p)\le\log K$，有 $0\le C\le1$。对概率
\begin{equation}
 \frac{\partial C}{\partial p_k}
 =\frac{\log p_k+1}{\log K}.
\end{equation}
通过 softmax logit $a_j$，
\begin{equation}
 \frac{\partial C}{\partial a_j}
 =\sum_k\frac{\partial C}{\partial p_k}p_k(\delta_{kj}-p_j).
\end{equation}
高确定度只表示分布尖锐，不表示 calibration；错误 one-hot 同样有 $C=1$。

\section{双 tick 选择的次梯度}
损失
\begin{equation}
 \Loss=\frac12(L_{t_1}+L_{t_2}),
 \quad t_1=\argmin_tL_t,
 \quad t_2=\argmax_tC_t.
\end{equation}
在唯一极值且索引不变的邻域，
\begin{equation}
 \nabla\Loss=\frac12(\nabla L_{t_1}+
abla L_{t_2}).
\end{equation}
当两个 tick 值相等时索引发生跳变，目标不可微，可取活跃极值集合梯度的凸包作为 Clarke 次梯度。实践的 argmin/argmax 选择等价于选择其中一个子梯度。

\section{BPTT 的时间 Jacobian}
递归状态
\begin{equation}
 h_{t+1}=F_\theta(h_t,o_t).
\end{equation}
最终损失对早期状态
\begin{equation}
 \frac{\partial\Loss}{\partial h_t}
 =\prod_{s=t}^{T-1}J_{F,h}(h_s)
 \frac{\partial\Loss}{\partial h_T}+\cdots.
\end{equation}
若 Jacobian 谱半径长期大于 1，梯度爆炸；小于 1，梯度消失。CTM 虽在内部引入多 tick，仍需处理普通递归网络的稳定性问题，通常依靠归一化、残差和截断 tick 数。

\section{自适应停机的风险--计算目标}
停机时间
\begin{equation}
 T_\tau=\min\{t:C_t\ge\tau\}\wedge T_{\max}.
\end{equation}
可把计算惩罚纳入训练
\begin{equation}
 J(\theta)=\E[L_{T_\tau}+\lambda T_\tau].
\end{equation}
$T_\tau$ 离散不可微，可用每 tick 停机概率 $h_t=\sigma(a_t)$ 构造可微期望
\begin{equation}
 \Pr(T=t)=h_t\prod_{s<t}(1-h_s),
\end{equation}
再优化
\begin{equation}
 \sum_t\Pr(T=t)(L_t+\lambda t).
\end{equation}
论文的阈值停机更简单，但阈值必须在验证集上校准。
\clearpage
\section{CTM 的双时间尺度}
外部数据步记为 $t$，内部思考 tick 记为 $s=1,\ldots,S$。对固定输入 $x_t$，模型在内部时间上递推
\begin{equation}
 h^{(s+1)}=F_\theta(h^{(s)},x_t,a^{(s)}).
\end{equation}
输出可在任意 tick 读取
\begin{equation}
 y^{(s)}=G_\theta(h^{(s)}).
\end{equation}
因此计算图不是普通“一层一表示”，而是同一输入上形成长度 $S$ 的内部动力系统。

\section{每个神经元的私人时间模型}
设第 $i$ 个神经元接收过去 $M$ 个预激活
\begin{equation}
 z_i^{(s-M+1:s)}.
\end{equation}
私人时间模型可写为
\begin{equation}
 h_i^{(s)}=\phi_i\left(
 \sum_{\tau=0}^{M-1}w_{i\tau}z_i^{(s-\tau)}+b_i
 \right).
\end{equation}
不同神经元拥有不同时间核 $w_i$，相当于 depthwise temporal convolution。若再用 MLP，
\begin{equation}
 h_i^{(s)}=g_i(z_i^{(s-M+1:s)}),
\end{equation}
则每个神经元可以学习非线性时间模式。

\section{时间核的频率响应}
线性部分
\begin{equation}
 y_i^{(s)}=\sum_{\tau=0}^{M-1}w_{i\tau}z_i^{(s-\tau)}
\end{equation}
的离散频率响应
\begin{equation}
 H_i(e^{i\omega})=
 \sum_{\tau=0}^{M-1}w_{i\tau}e^{-i\omega\tau}.
\end{equation}
某些神经元可学低通核，整合长期趋势；另一些可学高通或带通核，响应变化和振荡。CTM 的异质时间处理在数学上表现为不同 $H_i(\omega)$。

\section{同步矩阵是 Gram 矩阵}
设每个神经元在最近窗口的活动向量
\begin{equation}
 a_i=(h_i^{(s-M+1)},\ldots,h_i^{(s)})^\top\in\R^M.
\end{equation}
同步可定义为归一化内积
\begin{equation}
 S_{ij}=\frac{a_i^\top a_j}
 {\|a_i\|\|a_j\|+\epsilon}.
\end{equation}
若先归一化 $\widetilde a_i=a_i/\|a_i\|$，则
\begin{equation}
 S=\widetilde A^\top\widetilde A.
\end{equation}
所以 $S\succeq0$，对角线为 1，秩不超过窗口长度 $M$。同步表示并非任意矩阵，而受 Gram 几何约束。

\section{指数衰减同步的递推}
对乘积 $h_i^{(s)}h_j^{(s)}$ 做指数移动平均：
\begin{equation}
 C_{ij}^{(s)}=\lambda C_{ij}^{(s-1)}
 +(1-\lambda)h_i^{(s)}h_j^{(s)}.
\end{equation}
展开
\begin{equation}
 C_{ij}^{(s)}=(1-\lambda)
 \sum_{k=0}^{s-1}\lambda^k
 h_i^{(s-k)}h_j^{(s-k)}+
 \lambda^sC_{ij}^{(0)}.
\end{equation}
有效记忆长度约
\begin{equation}
 \tau_{\rm eff}\approx\frac1{1-\lambda}.
\end{equation}
例如 $\lambda=0.99$，约整合最近 100 个 tick。

\section{归一化相关系数的递推}
同时维护
\begin{equation}
 V_i^{(s)}=\lambda V_i^{(s-1)}
 +(1-\lambda)(h_i^{(s)})^2.
\end{equation}
归一化同步
\begin{equation}
 S_{ij}^{(s)}=\frac{C_{ij}^{(s)}}
 {\sqrt{V_i^{(s)}V_j^{(s)}}+\epsilon}.
\end{equation}
若所有活动乘同一尺度 $c$，分子乘 $c^2$，分母也乘 $c^2$，同步近似不变。因此它主要编码共同变化方向而非绝对幅值。

\section{随机 pair subsampling 的无偏估计}
完整同步有 $n(n-1)/2$ 对，成本 $O(n^2)$. 若均匀采样 $K$ 对集合 $\mathcal P$，估计平均同步损失
\begin{equation}
 \widehat L=\frac{M}{K}
 \sum_{(i,j)\in\mathcal P}\ell(S_{ij}),
 \qquad M=\frac{n(n-1)}2.
\end{equation}
若无放回均匀采样，
\begin{equation}
 \E[\widehat L]=
 \sum_{i<j}\ell(S_{ij})=L.
\end{equation}
方差随 $K$ 增大而下降。若采样概率不均匀 $p_{ij}$，需 Horvitz--Thompson 权重
\begin{equation}
 \widehat L=\frac1K\sum_{(i,j)\sim p}
 \frac{\ell(S_{ij})}{p_{ij}}.
\end{equation}

\section{同步表示驱动注意力的维度链}
选取 $K$ 个神经元对同步值组成
\begin{equation}
 s^{(s)}\in\R^K.
\end{equation}
线性投影为查询
\begin{equation}
 q^{(s)}=W_qs^{(s)}\in\R^{d_a}.
\end{equation}
外部视觉或序列 token 为
\begin{equation}
 X\in\R^{T\times d_x},
\end{equation}
键值
\begin{equation}
 K=XW_K\in\R^{T\times d_a},
 \quad V=XW_V\in\R^{T\times d_v}.
\end{equation}
注意力
\begin{equation}
 a^{(s)}=\softmax\left(
 \frac{q^{(s)}K^\top}{\sqrt{d_a}}
 \right)V\in\R^{d_v}.
\end{equation}
这样内部神经同步决定下一 tick 从外部输入读取什么。

\section{熵确定度的边界与梯度}
分类概率 $p\in\Delta^{C-1}$，熵
\begin{equation}
 H(p)=-\sum_{c=1}^{C}p_c\log p_c.
\end{equation}
归一化确定度
\begin{equation}
 \kappa=1-\frac{H(p)}{\log C}.
\end{equation}
均匀分布时 $H=\log C,\kappa=0$；one-hot 时 $H=0,\kappa=1$。对 logit $z_j$，利用 softmax Jacobian可得
\begin{equation}
 \frac{\partial H}{\partial z_j}
 =-p_j(\log p_j+H).
\end{equation}
确定度梯度为其相反数除以 $\log C$。

\section{双 tick 损失与不可导选择}
若训练选择最小损失 tick
\begin{equation}
 s_{\min}=\argmin_s\ell^{(s)},
\end{equation}
最大确定度 tick
\begin{equation}
 s_{\max}=\argmax_s\kappa^{(s)},
\end{equation}
目标
\begin{equation}
 \Loss=\ell^{(s_{\min})}+\ell^{(s_{\max})}.
\end{equation}
argmin/argmax 对参数是分段常数，反向只流过被选 tick。选择边界处梯度不连续。软替代可用
\begin{equation}
 w_s=\frac{e^{-\ell^{(s)}/\tau}}
 {\sum_re^{-\ell^{(r)}/\tau}},
 \qquad
 \Loss_{\rm soft}=\sum_sw_s\ell^{(s)},
\end{equation}
温度 $\tau\to0$ 时逼近最小值。

\section{BPTT 时间 Jacobian}
内部递推 $h^{(s+1)}=F(h^{(s)})$，最终损失对早期状态
\begin{equation}
 \frac{\partial\Loss}{\partial h^{(s)}}
 =\left(\prod_{r=s}^{S-1}J_r^\top\right)
 \frac{\partial\Loss}{\partial h^{(S)}}.
\end{equation}
若 $\|J_r\|_2\le\rho$，
\begin{equation}
 \left\|\frac{\partial\Loss}{\partial h^{(s)}}\right\|
 \le\rho^{S-s}
 \left\|\frac{\partial\Loss}{\partial h^{(S)}}\right\|.
\end{equation}
$\rho<1$ 梯度衰减，$\rho>1$ 爆炸。CTM 的内部 tick 增加了自适应计算，也带来与 RNN 相同的长时梯度问题。

\section{自适应停机的风险--计算目标}
设每个 tick 的停机概率 $p_s$，首次停机分布
\begin{equation}
 q_s=p_s\prod_{r<s}(1-p_r).
\end{equation}
期望任务损失
\begin{equation}
 \E_q[\ell]=\sum_sq_s\ell^{(s)},
\end{equation}
期望计算
\begin{equation}
 \E_q[s]=\sum_ssq_s.
\end{equation}
联合目标
\begin{equation}
 \Loss_{\rm ACT}=\sum_sq_s\ell^{(s)}+
 \lambda\sum_ssq_s.
\end{equation}
$\lambda$ 控制准确率与思考时间。仅用置信度阈值停机可能在过度自信样本上过早停止，因此需要校准。

\section{三个神经元的同步数值例子}
最近三 tick 活动
\begin{equation}
 a_1=(1,2,3),\quad
 a_2=(2,4,6),\quad
 a_3=(3,2,1).
\end{equation}
$a_2=2a_1$，余弦同步
\begin{equation}
 S_{12}=1.
\end{equation}
而
\begin{equation}
 S_{13}=\frac{1\cdot3+2\cdot2+3\cdot1}
 {\sqrt{14}\sqrt{14}}=\frac{10}{14}\approx0.714.
\end{equation}
同步表示能区分“同方向不同尺度”和“相反时间趋势”。
\end{document}
